\documentclass[12pt]{article}
\usepackage{amsmath,amsfonts,fullpage,amsthm,graphicx}

\newtheorem*{theorem}{Theorem}
\newtheorem*{fact}{Fact}
\newtheorem*{goal}{Goal}
\newtheorem*{idea}{Idea}
\newtheorem*{defn}{Definition}
\newtheorem*{ex}{Example}

%Style section
%\setlength{\textheight}{10in}
%\setlength{\textwidth}{7in}
%\hoffset=-0.75in
 \pagestyle{plain}
% \setlength{\topmargin}{0in}
% \setlength{\headsep}{0in}
% \setlength{\oddsidemargin}{0in}
% \setlength{\evensidemargin}{0in}

\begin{document}


\pagestyle{empty} %use this to get rid of page numbers

\noindent
{Math 166 \hskip 1.5 truein  Names:\ \hrulefill}

\noindent
%\vskip 0.3truein
%\bigskip
%\noindent
{Volumes using the disc and washer method}

%\noindent  \hskip 2.8 truein  Section Number:\ \hrulefill
\medskip

\begin{goal} \rm
Use integration to find the volume of various \emph{solids of revolution}.
\end{goal}
%\smallskip

\noindent {\bf Overview.} Let $f(x)$ be a positive, continuous function.  If the area underneath $f(x)$ from $x = a$ to $x = b$ is rotated around the $x$-axis, a $3$-dimensional object is formed, and the volume of that object is given by the formula,
\[
V = \pi \int_a^b f(x)^2 \, dx
\]
\begin{enumerate}
\item Explain why the area of a cross section at $x=a$ is given by $\pi f(a)^2$.  Then explain why the whole formula is valid. 
\vfill

\item Consider rotating the region under the function $f(x) = x$ from $x = 0$ to $x = 4$ about the $x$-axis.
\begin{enumerate}
\item Sketch a picture of the solid.  What type of solid is it?
\vfill
%\newpage
\item Use an integral to compute the volume of the solid.
\vfill
\item Use a known volume formula (you can Google one) to verify your answer from (b) is correct.
\vfill
\end{enumerate}
%\end{enumerate}

%In the following questions, first graph the function(s). Determine whether to integrate with respect to $x$ or $y$. State whether the cross-section is a disk or a washer. Then set up the integral(s). If you have time, evaluate the integral(s) to find the volume.
%
%\begin{enumerate}
%%\setcounter{enumi}{-1}

%\begin{enumerate}

\pagebreak
\item Consider the solid obtained by revolving the region enclosed by $y=32-2x$, $y=2x+4$, and $x=0$ about the $x$-axis.
\begin{enumerate}
\item Sketch the solid.
\vfill

%\newpage

\item What are the vertical cross-sections of this solid?
\vfill
\item How can we modify the volume formula at the top of the worksheet to find the volume of this solid?
\vfill

\end{enumerate}
\end{enumerate}
%\vspace{3in}
%
%\item Now find the volume of the solid obtained by revolving that same region about the $y$-axis. (Hint: You will need to use two integrals.)
%
%\vspace{3.5in}
%\pagebreak
%
%\item Find the volume of the solid obtained by revolving the region enclosed by $y=x^2$, $y=12-x$, and $x=-1$ about the line $y=15$. (Hint: Use your graph to figure out what your radius needs to be since we are revolving about the line $y=15$ instead of the $x$-axis.) 
%
%\vspace{3.75in}
%
%
%\item Sketch the hyperbola $y^2-x^2=1$ by plugging in $-2,1,0,1,2$ for $x$ and solving for $y$. (Hint: there will be two $y$ values for each $x$ value.) Find the volume of the solid obtained by revolving the region between the branches of the hyperbola for $-2\leq x\leq 2$ about the $x$-axis. Note the same solid would result if you revolved the region bounded by the top branch of the hyperbola, the lines $x=2$ and $x=-2$, and the $x$-axis, about the $x$-axis. (This solid is called a \emph{hyperboloid}.)
%
%
%%\newpage
%%
%%\item A region in the plane is rotated about the $x$-axis and the corresponding object has volume given by 
%%\[
%%\displaystyle \int_0^{12} \pi \left[144 - x^2\right] \, dx.
%%\]
%%\begin{enumerate}
%%\item Sketch an example of an object with this volume and cross-sections that are washers.  Mirroring the statements of problems 1-4, come up with a description of the region that is rotated about the $x$-axis.
%%\vfill
%%\item Do the same if the cross-sections are disks.
%%\vfill
%%\end{enumerate}
%
%\end{enumerate} 
%
\end{document}
